% Pursuing Stacks (cotes 134-2 to 134-8): archive cote 134-8, site page 12 (PDF page 13), typescript page 566.
% typescript. Draft compared: ps7.tex:555-613.
% AI-transcribed by Claude Opus 5.5 (claude-opus-5-5, adaptive thinking, effort high), Message Batches API, batch ps4 (msgbatch_012Aoc1NuskATPsyBeKSYK7J), 2026-11-12; not yet checked by eye.
\scanpage{134-8/12}
\mspage{566}
functor (more accurately, the composition of the latter with the
localization functor $\mathrm{Ch}_{\bullet}(\mathrm{Ab})\longrightarrow \mathrm{D}_{\bullet}(\mathrm{Ab}) = \mathrm{Hotab}$ ). T\over{h}{h}e equivalence of
the two definitions\inserted{(of abelianization)} is mentioned on the same p. 569, somewhat as a
chore I did'nt really feel then to dive into. \xout{The} Besides the ``block''
against homology, the picture was being obscured too by the ``computati-%
onal'' idea I kept in mind, of expressing $\mathrm{Lp}_{\hand{\unsure{!}}}^{ab}(F)$ for F in $A^{\wedge}_{ab}$ in
terms of an ``integrator'' for A i.e. as $F *_{\underline{Z}} L^{B}_{\bullet}$ - whereas I should
have known best myself that for establishing formal properties relating
various derived functors, the particular \xout{computational} approach\hand{es} used
for``co\inserted{m}puting'' them more or less eleg\xout{\unsure{g}}antly \unsure{are} wholly irrelevant ...

O\over{n}{n}e teaching I am getting out of all this, is that when expres-%
sing abelianization, or presumably any other kind of notion or operati-%
on of significance for homotopy types, one should be careful, for
any modelizer one chooses to work in, to dig out the description which
fits smoothly those particular models. Clearly, when working with
semisimplicial models, the description via tautological abelianization
$\mathrm{Wh}_{\hand{\Delta}}$ and using Kan-Dold-Puppe is the best. When working within the
modelizer (Cat) though, making the detour through $\hand{\Delta}^{\wedge}$ is awkward and
makes us \inserted{just} miss \deleted{the understanding of} the relevant facts. Once we got th\over{e}{is}
\deleted{understanding}, it gets clear too what to do when $\hand{\Delta}$ is replaced by any
other object A in (Cat) \xout{(which needs not even be a pseudotest category)}
when taking models in $A^{\wedge}$ (never minding even whether $A^{\wedge}$ is indeed a
``modelizer'' ): namely, apply $\mathrm{Wh}_{A}$ , and take total homology !

This brings to my mind another example, very similar indeed.
Some time after I got across Thomason's nice paper, showing that
(Cat) is a closed model category (see comments in section 87), I got
from Tim Porter another \xout{re} reprint of Thomason's, where he gives the
description of ``homotopy colimits'' (or ``integration'', as I call it)
within the modelizer (Cat),\xout{(compare} in terms of the total category
associated to a fibered category (compare section 69). There he grinds
through a tedious, highly technical proof, whereas the direct proof
when describing ``integration'' as the left adjoint functor to the
tautological ``inverse image'' functor, is more or less tautological
too. T\over{h}{h}e reason for this awkwardness is again that, rather than
being content to work with the models as they are, Thomason refers to
the Bousfield-Kan description of colimits in the semisimplicial set-up, \inserted{which he takes as his definition for colimits.}
I suspect that (Cat) is the one modelizer most ideally suited for expres%
-sing the ``integration'' operation, and that the Bousfield-Kan descripti-%
on is just the obvious, not-quite-as-simple one which can be deduced
\margin{12}
